\chapter{Journal de travail}
Cette annexe traite les différentes étapes réalisée au cours du projet. Cela permet d'avoir un court suivis en plus du Gantt. Le Gantt reste tout de même le moyen d'avoir un suivis plus précis.


\subsection*{16 février 2026}

\begin{itemize}
    \item Lecture du rapport de Thomas Freyche.
    \item Tentative de mise en place du rapport sous VS Code en latex(Echec).
\end{itemize}

\subsection*{18 février 2026}

\begin{itemize}
    \item Lecture du rapport de Thomas Freyche.
    \item Entretien avec M.Girardin.
\end{itemize}

\subsection*{19 février 2026}

\begin{itemize}
    \item Lecture du rapport de Thomas Freyche.
    \item Tentative de mise en place du rapport sous VS Code en latex(Echec).
\end{itemize}

\subsection*{20 février 2026}

\begin{itemize}
    \item Lecture du rapport de Thomas Freyche.
    \item Tentative de mise en place du rapport sous VS Code en latex(Réussite).
\end{itemize}

\subsection*{23 février 2026}

\begin{itemize}
    \item Recherche d'un protocole pour réaliser une première mise en sérvice.
    \item Documentation de la mise en service et des étapes réaliser pour y réussir.
\end{itemize}

\subsection*{25 février 2026}

\begin{itemize}
    \item Première essaie de la mise en service
    \item Entretien avec M.Girardin
\end{itemize}

\subsection*{26 février 2026}

\begin{itemize}
    \item Création d'un Gantt pour un bons suivis du TB
\end{itemize}

\subsection*{27 février 2026}

\begin{itemize}
    \item Suite de la création d'un Gantt pour un bons suivis du TB
    \item Courte mise en place de la structure de la documentation
\end{itemize}

\subsection*{2 mars 2026}
\begin{itemize}
    \item Discussion avec Thomas Freyche
    \item Lecture du code (Analyse des librairies)
    \item Documentation du code dans le rapport
    \item Entretien avec M.Girardin
\end{itemize}

\subsection*{5 mars 2026}
\begin{itemize}
    \item Écriture du rapport (Ajout de section)
\end{itemize}

\subsection*{5 mars 2026}
\begin{itemize}
    \item Suite de la lecture du code
    \item Documentation du code dans le rapport
    \item Ajout de nouvelle section dans le rapport
    \item Écriture du rapport (Ajout de sections)
\end{itemize}

\subsection*{6 mars 2026}
\begin{itemize}
    \item Suite de la lecture du code
    \item Documentation du code dans le rapport
    \item Réorganisation du rapport
\end{itemize}

\subsection*{7 mars 2026}
\begin{itemize}
    \item Suite de la lecture du code
    \item Documentation du code dans le rapport
    \item Réorganisation du rapport
\end{itemize}

\subsection*{8 mars 2026}
\begin{itemize}
    \item Suite de la lecture du code
    \item Documentation du code dans le rapport
    \item Réorganisation du rapport
\end{itemize}

\subsection*{16 mars 2026}
\begin{itemize}
    \item Suite de la lecture du code
    \item Documentation du code dans le rapport
\end{itemize}

\subsection*{17 mars 2026}
\begin{itemize}
    \item Recherche d'information sur l'ESP32
    \item Suite de l'analyse du code
\end{itemize}

\subsection*{22 mars 2026}
\begin{itemize}
    \item Recherche d'information sur l'ESP32
\end{itemize}

\subsection*{23 mars 2026}
\begin{itemize}
    \item Discussion avec M. Herzog
    \item Lecture et documation du code
    \item Création de la machine d'état
\end{itemize}

\subsection*{24 mars 2026}
\begin{itemize}
    \item Fin de la lecture et documation du code
\end{itemize}

\subsection*{26 mars 2026}
\begin{itemize}
    \item Embellsiment et amélioration sur le code de sustentation
    \item Documentation et planification
    \item Suite de la machine d'état
\end{itemize}

\subsection*{27 mars 2026}
\begin{itemize}
    \item Embellsiment et amélioration sur le code de sustentation
    \item Documentation
\end{itemize}

\subsection*{30 mars 2026}
\begin{itemize}
    \item Création des fichiers de mesures de temps
    \item Mesure des temps
    \item Capture d'oscillo
\end{itemize}

\subsection*{31 mars 2026}
\begin{itemize}
    \item Finalisation de la machine d'état
\end{itemize}

\subsection*{2 avril 2026}
\begin{itemize}
    \item Finalisation des mesures de temps
    \item Début de l'écriture en C de la machine d'état
\end{itemize}

\subsection*{12 avril 2026}
\begin{itemize}
    \item Finalisation de la machine d'état
\end{itemize}

\subsection*{13 avril 2026}
\begin{itemize}
    \item Test de la machine d'état
    \item Mesure des positions
    \item Recherche d'information concernant le capteur
    \item Récupération de cales
\end{itemize}

\subsection*{14 avril 2026}
\begin{itemize}
    \item Documentation
    \item Première mesure
    \item Recherche de cales
\end{itemize}

\subsection*{15 avril 2026}
\begin{itemize}
    \item Mesure de positons (Tension et distance)
    \item Vérification des composants du conditionneur
    \item Recherche d'information sur le calibrage du capteur
    \item Implémentation d'un filtre pour les positions
    \item Limer des cales
\end{itemize}

\subsection*{16 avril 2026}
\begin{itemize}
    \item Documentation
\end{itemize}

\subsection*{17 avril 2026}
\begin{itemize}
    \item Documentation et planification
\end{itemize}

\subsection*{20 avril 2026}
Voici le travail réalisé :

\begin{itemize}
    \item Recalibration du capteur
    \item Mesure de positions de l'inducteur 3
    \item Essaie d'une calibration pour vérifier si le capteur de position peut donner une valeur plus grande que 3V.
    \item Analyse de la lévitation suite à la recalibration.
\end{itemize}

Voici les choses observées :

\begin{itemize}
    \item On a essayé de mettre une calle sur l'inducteur 4 pour élever un peu l'inducteur 3. Enf aisant cela, nous avons recalibré le capteur On souhaitait voir si le capteur de position peut déapsser les 3V ou alors est-ce qu'il se bloque à 3V lorsque la position à une valeur plus grande que l'entrefer. Le résultat obtenu montre que le capteur se bloque à 3V meme si la distance est plus grande que 3mm.
    \item Une petite anomalie détectée durant une mesure de position a étét que le sginal ne restait pas droit. Les conditions résultantes de cela était qu'aucun inducteur n'était connecté. Lors de l'enclenchement du switch de précharge, un signal en creux à 25kHz apparaissait.
    \item Suite à la calibration du capteur, une lévitation fu réaliser avec le code le plus récent et celui du travail précédent. Le code du travail précédent donnait le meme résultat que celui de départ. Le code le plus récent à cette date, donnait un blocage des inducteurs de droite à un entrefer de 0 mm et ceux de gauche un entrefer de 3mm. Lorsque la régulation est réalisée que sur l'avant, on constate que l'inducteur 1 ne se lève pas beaucoup. Il a été décidé de mesurer les courants pour vérifier la force max possible à atteindre lorsqu'on applique une force en plus de celle de pesanteur.
\end{itemize}

\subsection*{21 avril 2026}
\begin{itemize}
    \item Documentation
    \item Création d'un tableau de mesure général
\end{itemize}

\subsection*{22 avril 2026}
Travail réalisé :

\begin{itemize}
    \item Préparations du tableau de mesure
    \item Mesure des alimentations
\end{itemize}

Observation :
\begin{itemize}
    \item Le régulateur 24V a besoin d'une pression physique pour s'enclencher. Ce n'est pas normal. A cause de celui-la, l'alimentation générale n'est pas fonctionelle. Un arc électrique se crée sur une de pin et une autre. En rebrassant les connecteurs, le régulateur fonctionne une nouvelle fois.
\end{itemize}

\subsection*{23 avril 2026}
Travail réalisé :

\begin{itemize}
    \item Finalisation du tableau
    \item Première mesure des capteurs de courants
    \item Mesure général des capteurs de courants
    \item Première mesure général des capteurs de position
\end{itemize}

Observation :
\begin{itemize}
    \item Le régulateur 24V semble aller mieux. Cela devait surement être la pin mentionnée.
\end{itemize}

\subsection*{24 avril 2026}
\begin{itemize}
    \item Mesure général du capteur de positions
    \item Suite des mesure des capteurs de courants
    \item Mesure des capteurs de positions
    \item Documentation
\end{itemize}

\subsection*{25 avril 2026}
\begin{itemize}
    \item Documentation
\end{itemize}

\subsection*{26 avril 2026}
\begin{itemize}
    \item Documentation et planification
\end{itemize}